\documentclass[12pt,oneside]{book}

% -------------------------------------------------
% PAQUETES
% -------------------------------------------------
\usepackage[spanish,es-nodecimaldot]{babel}
\usepackage[utf8]{inputenc}
\usepackage[T1]{fontenc}
\usepackage{lmodern}
\usepackage[a4paper,margin=2.5cm]{geometry}
\usepackage{microtype}
\usepackage{setspace}
\usepackage{graphicx}
\usepackage{xcolor}
\usepackage{fancyhdr}
\usepackage{titlesec}
\usepackage{hyperref}
\usepackage{bookmark}
\usepackage{tocloft}
\usepackage{tcolorbox}
\tcbuselibrary{theorems, skins, breakable}
\usepackage{amsmath,amssymb}
\usepackage{enumitem}
\usepackage{csquotes}

% -------------------------------------------------
% CONFIGURACIÓN VISUAL  (paleta cósmica)
% -------------------------------------------------
\definecolor{azulCosmos}{HTML}{1B3A6B}
\definecolor{azulCielo}{HTML}{2980B9}
\definecolor{doradoEstrella}{HTML}{D4A017}
\definecolor{grisClaro}{HTML}{F4F6F8}
\definecolor{grisTexto}{HTML}{2E2E2E}
\definecolor{grisSec}{HTML}{6C757D}
\definecolor{colorConcepto}{HTML}{EBF5FB}
\definecolor{colorActividad}{HTML}{EAFAF1}
\definecolor{colorIdea}{HTML}{FEF9E7}
\definecolor{colorNota}{HTML}{F2F3F4}
\definecolor{colorCompromiso}{HTML}{EAF0FB}
\definecolor{colorPeligro}{HTML}{FDEDEC}
\definecolor{rojoPeligro}{HTML}{C0392B}

\hypersetup{
    colorlinks=true,
    linkcolor=azulCosmos,
    urlcolor=azulCielo,
    citecolor=azulCosmos,
    pdftitle={Consulta — Actividad 2.1: Normas de Uso del Telescopio},
    pdfauthor={Semillero de Astronomía — IEEOH}
}

\setstretch{1.12}
\setlength{\parskip}{0.4em}
\setlength{\parindent}{0pt}
\setcounter{secnumdepth}{3}
\setcounter{tocdepth}{2}
\setlength{\headheight}{26pt}

% -------------------------------------------------
% ESTILOS DE TÍTULOS
% -------------------------------------------------
\titleformat{\chapter}[display]
  {\bfseries\Huge\color{azulCosmos}}
  {\filleft\Large\chaptername\ \thechapter}
  {1ex}
  {\titlerule\vspace{1ex}\filright}
  [\vspace{1ex}\titlerule]

\titleformat{\section}
  {\bfseries\Large\color{azulCosmos}}
  {\thesection}{0.7em}{}

\titleformat{\subsection}
  {\bfseries\large\color{grisTexto}}
  {\thesubsection}{0.6em}{}

\titleformat{\subsubsection}
  {\bfseries\normalsize\color{grisTexto}}
  {\thesubsubsection}{0.5em}{}

% -------------------------------------------------
% ENCABEZADOS Y PIES
% -------------------------------------------------
\pagestyle{fancy}
\fancyhf{}
\fancyhead[L]{\textcolor{grisSec}{\nouppercase{\leftmark}}}
\fancyhead[R]{\textcolor{grisSec}{Semillero de Astronomía — IEEOH}}
\fancyfoot[C]{\textcolor{grisSec}{\thepage}}
\renewcommand{\headrulewidth}{0.4pt}
\renewcommand{\headrule}{%
  \hbox to\headwidth{\color{azulCosmos}\leaders\hrule height \headrulewidth\hfill}}

% -------------------------------------------------
% ÍNDICE
% -------------------------------------------------
\renewcommand{\cftchapfont}{\bfseries\color{azulCosmos}}
\renewcommand{\cftsecfont}{\color{grisTexto}}
\renewcommand{\cftchappagefont}{\bfseries\color{azulCosmos}}
\renewcommand{\cftsecpagefont}{\color{grisTexto}}

% -------------------------------------------------
% ENTORNOS TCOLORBOX
% -------------------------------------------------

\newtcolorbox{destacado}{
  colback=grisClaro,
  colframe=azulCosmos,
  boxrule=0.8pt,
  arc=2mm,
  left=3mm, right=3mm, top=1.5mm, bottom=1.5mm
}

\newtcbtheorem[number within=chapter]{concepto}{Concepto}{
  enhanced,
  colback=colorConcepto,
  colframe=azulCosmos,
  coltitle=white,
  colbacktitle=azulCosmos,
  attach boxed title to top left={yshift=-2mm, xshift=4mm},
  boxed title style={enhanced, arc=1mm, boxrule=0pt, colframe=azulCosmos},
  fonttitle=\bfseries,
  breakable,
  arc=1.5mm,
  boxrule=0.8pt,
  left=4mm, right=4mm, top=5mm, bottom=2mm,
  separator sign none,
  description delimiters={(}{)}
}{con}

\newtcbtheorem[number within=chapter]{actividad}{Actividad}{
  enhanced,
  colback=colorActividad,
  colframe=azulCosmos,
  coltitle=white,
  colbacktitle=azulCosmos,
  attach boxed title to top left={yshift=-2mm, xshift=4mm},
  boxed title style={enhanced, arc=1mm, boxrule=0pt, colframe=azulCosmos},
  fonttitle=\bfseries,
  breakable,
  arc=1.5mm,
  boxrule=0.8pt,
  left=4mm, right=4mm, top=5mm, bottom=2mm,
  separator sign none,
  description delimiters={(}{)}
}{act}

\newtcbtheorem[number within=chapter]{idea}{Idea}{
  enhanced,
  colback=colorIdea,
  colframe=doradoEstrella,
  coltitle=white,
  colbacktitle=doradoEstrella,
  attach boxed title to top left={yshift=-2mm, xshift=4mm},
  boxed title style={enhanced, arc=1mm, boxrule=0pt, colframe=doradoEstrella},
  fonttitle=\bfseries,
  breakable,
  arc=1.5mm,
  boxrule=0.8pt,
  left=4mm, right=4mm, top=5mm, bottom=2mm,
  separator sign none,
  description delimiters={(}{)}
}{idea}

\newtcbtheorem[number within=chapter]{nota}{Nota}{
  enhanced,
  colback=colorNota,
  colframe=grisSec,
  coltitle=white,
  colbacktitle=grisSec,
  attach boxed title to top left={yshift=-2mm, xshift=4mm},
  boxed title style={enhanced, arc=1mm, boxrule=0pt, colframe=grisSec},
  fonttitle=\bfseries,
  breakable,
  arc=1.5mm,
  boxrule=0.6pt,
  left=4mm, right=4mm, top=5mm, bottom=2mm,
  separator sign none,
  description delimiters={(}{)}
}{nota}

\newtcbtheorem[number within=chapter]{compromiso}{Compromiso}{
  enhanced,
  colback=colorCompromiso,
  colframe=azulCielo,
  coltitle=white,
  colbacktitle=azulCielo,
  attach boxed title to top left={yshift=-2mm, xshift=4mm},
  boxed title style={enhanced, arc=1mm, boxrule=0pt, colframe=azulCielo},
  fonttitle=\bfseries,
  breakable,
  arc=1.5mm,
  boxrule=0.8pt,
  left=4mm, right=4mm, top=5mm, bottom=2mm,
  separator sign none,
  description delimiters={(}{)}
}{comp}

\newtcbtheorem[number within=chapter]{regla}{Regla}{
  enhanced,
  colback=grisClaro,
  colframe=azulCosmos,
  coltitle=white,
  colbacktitle=azulCosmos,
  attach boxed title to top left={yshift=-2mm, xshift=4mm},
  boxed title style={enhanced, arc=1mm, boxrule=0pt, colframe=azulCosmos},
  fonttitle=\bfseries,
  breakable,
  arc=1.5mm,
  boxrule=0.8pt,
  left=4mm, right=4mm, top=5mm, bottom=2mm,
  separator sign none,
  description delimiters={(}{)}
}{reg}

% Caja especial de advertencia / peligro
\newtcbtheorem[number within=chapter]{peligro}{Escenario hipot\'etico}{
  enhanced,
  colback=colorPeligro,
  colframe=rojoPeligro,
  coltitle=white,
  colbacktitle=rojoPeligro,
  attach boxed title to top left={yshift=-2mm, xshift=4mm},
  boxed title style={enhanced, arc=1mm, boxrule=0pt, colframe=rojoPeligro},
  fonttitle=\bfseries,
  breakable,
  arc=1.5mm,
  boxrule=0.8pt,
  left=4mm, right=4mm, top=5mm, bottom=2mm,
  separator sign none,
  description delimiters={(}{)}
}{pel}

% -------------------------------------------------
% DATOS DEL DOCUMENTO
% -------------------------------------------------
\newcommand{\NombreSemillero}{Semillero de Astronomía}
\newcommand{\Institucion}{Institución Educativa Enrique Olaya Herrera}
\newcommand{\Ciudad}{Medellín, Colombia}
\newcommand{\AsesorPrincipal}{Daniel Soto Villada}
\newcommand{\CoAsesor}{Gerardo Ríos}
\newcommand{\TituloDocumento}{Consulta — Actividad~2.1\\[0.4em]Normas de Uso\\[0.2em]del Telescopio}

% -------------------------------------------------
% DOCUMENTO
% -------------------------------------------------
\begin{document}
\hypersetup{pageanchor=false}

% ---------- PORTADA ----------
\begin{titlepage}
    \centering
    \vspace*{1.5cm}

    {\Large\textcolor{grisSec}{\Institucion}\par}
    \vspace{0.2cm}
    {\large\textcolor{grisSec}{\Ciudad}\par}

    \vspace{2.0cm}

    {\Huge\bfseries\textcolor{azulCosmos}{\TituloDocumento}\par}
    \vspace{0.5cm}
    {\Large\textcolor{grisSec}{Consultas Individuales — Segundo Encuentro}\par}

    \vspace{2.4cm}

    \begin{tcolorbox}[
        colback=grisClaro,
        colframe=azulCosmos,
        boxrule=0.9pt,
        arc=2mm,
        width=0.82\textwidth
    ]
        \centering
        {\large\bfseries\textcolor{azulCosmos}{Integrantes}\par}
        \vspace{0.3cm}
        {\normalsize Miguel \quad$\cdot$\quad Thomas \quad$\cdot$\quad Kevin\par}
        \vspace{0.1cm}
        {\normalsize Isabella \quad$\cdot$\quad Esteban \quad$\cdot$\quad Samuel Angulo\par}
        \vspace{0.4cm}
        {\large\bfseries\textcolor{azulCosmos}{Asesores}\par}
        \vspace{0.2cm}
        {\large \AsesorPrincipal\ \textcolor{grisSec}{(Director Académico)}\par}
        \vspace{0.1cm}
        {\large \CoAsesor\ \textcolor{grisSec}{(Director de Gestión)}\par}
    \end{tcolorbox}

    \vspace{1.6cm}

    \begin{tcolorbox}[
        colback=colorIdea,
        colframe=doradoEstrella,
        boxrule=0.7pt,
        arc=2mm,
        width=0.82\textwidth
    ]
        \centering
        \itshape\large\textcolor{grisTexto}{%
        ``Un instrumento bien cuidado es un instrumento que siempre estará listo
        para revelar los secretos del universo.''}
    \end{tcolorbox}

    \vfill
    {\large\textcolor{grisSec}{\today}\par}
    \vspace{0.5cm}
\end{titlepage}

% ---------- PÁGINAS PRELIMINARES ----------
\frontmatter
\tableofcontents
\clearpage

% ---------- CONTENIDO ----------
\clearpage
\hypersetup{pageanchor=true}
\mainmatter

% ================================================================
%  CAPÍTULO 1 — Normas de Uso del Telescopio
% ================================================================
\chapter{Normas de Uso del Telescopio}

\begin{destacado}
    \textbf{Actividad:} 2.1 \quad
    \textbf{Asignación:} Segundo Encuentro (29 de abril de 2026) \quad
    \textbf{Instrumento:} Telescopio \textit{Decouvir Maksutov-Cassegrain MC~60/700}
\end{destacado}

\medskip

En el segundo encuentro del semillero se asignó a cada integrante la investigación
de una norma específica de uso del telescopio. El objetivo es comprender no solo
\emph{qué} se debe hacer, sino \emph{por qué} esa norma existe, qué riesgo técnico
o de seguridad previene y qué consecuencias tendría ignorarla. Este documento recoge
los hallazgos de cada consulta, enriquecidos con ejemplos hipotéticos que ilustran
de manera concreta los riesgos asociados \cite{suiter1994,thompson2005,rutten1988}.

% ---------------------------------------------------------------
\section{Miguel — No apuntar al Sol sin filtro solar certificado}
% ---------------------------------------------------------------

\begin{concepto}{Concentración de radiación solar en un telescopio}{con:solar}
    Un telescopio no es simplemente un instrumento que \emph{magnifica} la imagen:
    es, ante todo, un colector de luz. Al apuntar al Sol, toda la radiación que
    entra por el objetivo se concentra en un punto focal de milímetros de diámetro.
    Para un telescopio con apertura de 60~mm, la energía solar que llega al ojo
    o al ocular es entre \textbf{600 y 3\,600 veces} mayor que la que incide sobre
    el ojo desnudo, dependiendo del modelo óptico \cite{pasachoff2000}. Esta
    concentración incluye radiación visible, ultravioleta e infrarroja.
\end{concepto}

\subsection{¿Por qué es importante esta norma?}

La retina humana carece de receptores de dolor. Esto significa que una lesión
grave puede ocurrir sin ningún aviso previo: \textbf{ni ardor, ni sensación de
calor, ni señal de alerta}. La concentración de radiación solar a través del
telescopio puede causar \emph{fotocoagulación retiniana} en milisegundos —~literalmente,
las proteínas de las células fotorreceptoras se desnaturalizan por calor, del mismo
modo que la clara de un huevo se coagula \cite{dobson2011}.

Existen dos tipos de filtros solares aptos para observación con telescopio:

\begin{itemize}[leftmargin=1.5em]
    \item \textbf{Filtros de vidrio óptico solar:} atenúan la luz en un factor
          de 100\,000:1 o más (densidad óptica $\geq 5$). Deben cumplir las normas
          ISO~12312-2 para observación directa del Sol \cite{iso12312}.
    \item \textbf{Película Baader AstroSolar:} lámina de polímero multicapa con
          densidad óptica 5.0 (versión visual), ampliamente utilizada en astronomía
          aficionada y recomendada por la comunidad científica \cite{baader2023}.
\end{itemize}

\begin{nota}{Filtros que NO son válidos}{nota:filtros-invalidos}
    Lentes de sol convencionales, negativos fotográficos velados, CDs, vidrios
    ahumados o cualquier improvisación casera \textbf{no son filtros solares seguros}.
    Pueden bloquear la luz visible pero transmiten radiación infrarroja invisible
    que igualmente destruye la retina. Solo los filtros certificados bajo estándar
    ISO garantizan protección completa \cite{pasachoff2000,iso12312}.
\end{nota}

\subsection{¿Qué riesgo se evita?}

\begin{itemize}[leftmargin=1.5em]
    \item \textbf{Ceguera permanente} por fotocoagulación retiniana: la lesión es
          irreversible porque las células de la retina no se regeneran.
    \item \textbf{Daño al ocular}: el calor concentrado puede fundir el cemento óptico
          de los elementos del ocular y deformar los barriles plásticos o los sellos de goma.
    \item \textbf{Daño al espejo o al corrector}: en telescopios tipo Maksutov, la
          lente correctora puede agrietarse por el gradiente térmico extremo.
\end{itemize}

\begin{peligro}{Una fracción de segundo que lo cambia todo}{pel:solar}
    \itshape
    Imaginemos que Miguel, durante una sesión de observación solar en el patio,
    retira el filtro Baader por \emph{``un segundo''} para ver si realmente es
    necesario. Al hacerlo, la energía solar concentrada incide directamente sobre
    su fóvea —~la zona de mayor agudeza visual de la retina. En ese instante,
    la temperatura local supera los 65\,°C: las células de los conos se desnaturalizan
    de forma irreversible. Miguel no siente dolor en el momento, pero al día siguiente
    nota una mancha oscura fija en el centro de su campo visual. El diagnóstico del
    oftalmólogo confirma una \emph{retinopatía solar} con escotoma central permanente.
    Nunca más podrá leer con normalidad ni ver detalles finos en astronomía.

    \medskip
    \normalfont
    Este escenario, documentado en numerosos casos clínicos tras eclipses solares
    sin protección adecuada, ilustra por qué esta norma no admite excepciones
    \cite{dobson2011,pasachoff2000}.
\end{peligro}

\begin{regla}{Regla de oro para la observación solar}{reg:filtro-solar}
    El filtro solar se coloca \textbf{siempre antes} de apuntar el telescopio al
    Sol, y se retira \textbf{siempre después} de apartar el telescopio del Sol.
    Nunca al revés. El filtro va en el extremo del objetivo (frente del tubo),
    no en el ocular.
\end{regla}

% ---------------------------------------------------------------
\section{Thomas — Aclimatación del telescopio antes de su uso}
% ---------------------------------------------------------------

\begin{concepto}{Equilibrio térmico y \emph{seeing} interno}{con:aclim}
    Todo material sólido, incluyendo los espejos y lentes de un telescopio, se
    expande o contrae al cambiar de temperatura. Cuando un telescopio pasa de un
    ambiente cálido al exterior nocturno, sus componentes tardan tiempo en alcanzar
    la temperatura ambiente. Mientras tanto, la diferencia de temperatura entre las
    superficies ópticas y el aire circundante genera \emph{corrientes de tubo}
    (turbulencias térmicas internas): capas de aire a distintas densidades que
    distorsionan la luz como lo hace el asfalto caliente en una carretera \cite{suiter1994}.
\end{concepto}

\subsection{¿Por qué es importante esta norma?}

En el telescopio \textit{Decouvir Maksutov-Cassegrain}, la lente correctora tipo
Maksutov tiene varios milímetros de espesor. Esta masa de vidrio actúa como un
\emph{acumulador térmico}: tarda considerablemente más en equilibrarse que los
espejos delgados de un Newtoniano equivalente. La aclimatación recomendada es de
\textbf{20 a 40 minutos} en exteriores antes de iniciar la sesión de observación,
dependiendo de la diferencia de temperatura \cite{tonkin2007}.

Durante la aclimatación, el telescopio debe permanecer \textbf{descubierto y en
posición vertical u horizontal}, sin apuntar al cielo todavía, para que el flujo de
aire natural ventile el interior del tubo de forma homogénea.

\begin{nota}{El efecto en Medellín}{nota:medellin-aclim}
    En Medellín la diferencia entre la temperatura interior de un aula (alrededor
    de 22\,°C) y el exterior en horas de madrugada puede alcanzar los 8–12\,°C.
    Aunque no es la mayor diferencia imaginable, sí es suficiente para generar
    corrientes de tubo visibles en observaciones planetarias de alto aumento.
\end{nota}

\subsection{¿Qué riesgo se evita?}

\begin{itemize}[leftmargin=1.5em]
    \item \textbf{Imágenes borrosas o inestables}: las corrientes de tubo hacen
          que estrellas puntuales parezcan danzar o deformarse.
    \item \textbf{Frustración del grupo}: sin saber la causa, los estudiantes
          podrían concluir erróneamente que el telescopio está dañado.
    \item \textbf{Pérdida de oportunidades astronómicas}: eventos únicos (ocultaciones,
          conjunciones, pasos de ISS) no esperan a que el equipo se aclimate.
\end{itemize}

\begin{peligro}{La noche de Saturno que nunca llegó a verse}{pel:aclim}
    \itshape
    El semillero planea una observación especial de Saturno en oposición, el mejor
    momento del año para verlo. Thomas saca el telescopio directamente del casillero
    del aula (23\,°C) a la terraza del colegio (13\,°C). Entusiasmados, apuntan de
    inmediato al planeta. A través del ocular de 10~mm ven algo redondo, pero los
    anillos parecen bordes difusos que tiemblan sin parar; nadie distingue la división
    de Cassini ni la sombra del planeta sobre los anillos. Frustrante.

    Treinta minutos después, cuando por fin el telescopio se aclimató, la imagen
    sería perfectamente estable y los anillos estarían claramente separados. Pero
    la sesión ya terminó porque era hora de regresar al aula. La oposición no se
    repetirá en casi un año.

    \medskip
    \normalfont
    La aclimatación no es un trámite: es el primer paso de toda sesión exitosa.
    \cite{suiter1994,tonkin2007}
\end{peligro}

% ---------------------------------------------------------------
\section{Kevin — Evitar el contacto directo con los elementos ópticos}
% ---------------------------------------------------------------

\begin{concepto}{Recubrimientos ópticos y contaminación de superficie}{con:optica}
    Los espejos y lentes de un telescopio poseen \emph{recubrimientos ópticos}:
    capas delgadas de aluminio, plata o dieléctricas de apenas decenas de nanómetros
    de espesor, aplicadas al vacío. Estos recubrimientos determinan la reflectividad
    y transmitancia del instrumento. Una huella dactilar deposita sobre la superficie
    una mezcla de aceites, sales, aminoácidos y agua con pH ligeramente ácido
    (alrededor de 5.5), que ataca química y mecánicamente estos recubrimientos de
    forma progresiva e irreversible si no se limpia con el procedimiento correcto
    \cite{smith2000}.
\end{concepto}

\subsection{¿Por qué es importante esta norma?}

Una sola huella dactilar sobre el espejo principal o la lente correctora puede:

\begin{itemize}[leftmargin=1.5em]
    \item \textbf{Reducir el contraste de la imagen} al dispersar la luz en múltiples
          direcciones (\emph{scattering}), haciendo que planetas brillantes generen
          halos difusos.
    \item \textbf{Degradar el recubrimiento permanentemente}: los ácidos grasos
          de la piel corroen el aluminio del espejo en horas si el ambiente es húmedo.
    \item \textbf{Provocar rayones} si se intenta limpiar la huella sin el procedimiento
          adecuado: un simple pañuelo de tela puede introducir partículas abrasivas que
          rayan el recubrimiento de manera irreparable.
\end{itemize}

La única limpieza válida de óptica astronómica es con:

\begin{enumerate}[leftmargin=1.5em]
    \item Aire comprimido limpio (perabuble o bombín fotográfico) para retirar polvo suelto.
    \item Hisopo de algodón y alcohol isopropílico $\geq 99\%$ en movimiento espiral
          desde el centro hacia afuera, con presión mínima \cite{thompson2005}.
\end{enumerate}

\begin{nota}{Una acción que parece inofensiva}{nota:huella}
    Tocar la óptica parece algo menor. Sin embargo, en la comunidad astronómica
    aficionada es uno de los daños más comunes y costosos. Reemplazar el recubrimiento
    de un espejo primario en un taller especializado puede costar entre \$50 y \$200
    USD, un valor que fácilmente supera el costo total del telescopio escolar.
\end{nota}

\subsection{¿Qué riesgo se evita?}

\begin{itemize}[leftmargin=1.5em]
    \item Pérdida de contraste y aparición de halos en las imágenes.
    \item Corrosión permanente del recubrimiento metálico.
    \item Rayones si se intenta limpiar incorrectamente.
    \item Gasto económico de reacondicionamiento o reemplazo del componente óptico.
\end{itemize}

\begin{peligro}{Cinco huellas, un espejo arruinado}{pel:huellas}
    \itshape
    Kevin tiene curiosidad por ver cómo se ve el espejo secundario del Maksutov
    desde adentro del tubo. Sin pensarlo demasiado, agarra el tubo por la lente
    correctora para asomarse. Sus cinco dedos dejan huellas sobre la lente.
    Al darse cuenta, toma un pañuelo de tela y frota para limpiar. El pañuelo
    arrastra una partícula de polvo microscópica y deja tres rayones de 2~cm sobre
    el recubrimiento antirreflejo.

    A partir de ese día, todas las observaciones de la Luna muestran un halo difuso
    alrededor del disco y las estrellas brillantes tienen halos de dispersión
    visibles. El recubrimiento dañado no se puede restaurar: la lente correctora
    deberá enviarse a reacondicionamiento o reemplazarse.

    \medskip
    \normalfont
    Lección: la óptica astronómica se mira, nunca se toca.
    \cite{smith2000,thompson2005}
\end{peligro}

% ---------------------------------------------------------------
\section{Isabella — Cubrir el objetivo y el ocular cuando el equipo no esté en uso}
% ---------------------------------------------------------------

\begin{concepto}{Protección de la óptica frente a polvo, humedad y hongos}{con:tapas}
    Las superficies ópticas de un telescopio son estructuralmente vulnerables a
    dos agentes ambientales: el \emph{polvo} y la \emph{humedad}. El polvo acumula
    partículas abrasivas que dañan el recubrimiento en la limpieza posterior.
    La humedad, especialmente en climas tropicales como el de Medellín, favorece
    la \emph{condensación} en superficies frías y —~en presencia de esporas
    presentes en el aire interior de edificios~— el crecimiento de \textbf{hongos
    ópticos}: colonias microscópicas que digieren literalmente el recubrimiento
    del vidrio en procesos lentos pero irreversibles \cite{rutten1988,cloudy2019}.
\end{concepto}

\subsection{¿Por qué es importante esta norma?}

Las tapas del objetivo y del ocular cumplen funciones distintas pero complementarias:

\begin{itemize}[leftmargin=1.5em]
    \item \textbf{Tapa del objetivo (parte frontal del tubo)}: protege la lente correctora
          o el espejo primario del polvo en suspensión, insectos, salpicaduras y luz
          parásita. En el Maksutov, también protege el interior del tubo herméticamente
          sellado, que no debe ventilarse con el aire húmedo del entorno.
    \item \textbf{Tapa del ocular (portaocular)}: cuando el telescopio está sin ocular
          colocado, el portaocular queda expuesto al espejo secundario; polvo y humedad
          pueden llegar directamente al sistema óptico interno.
\end{itemize}

\begin{nota}{El clima de Medellín y los hongos ópticos}{nota:hongos}
    Medellín tiene una humedad relativa promedio del 65–80\%. Los armarios y depósitos
    escolares pueden superar el 80\% de humedad relativa en temporada de lluvias.
    En estas condiciones, los \emph{hongos ópticos} (principalmente del género
    \textit{Penicillium} y \textit{Aspergillus}) pueden colonizar superficies
    ópticas no protegidas en semanas. Una vez establecidos, son prácticamente
    imposibles de eliminar sin dañar el recubrimiento \cite{rutten1988}.
\end{nota}

\subsection{¿Qué riesgo se evita?}

\begin{itemize}[leftmargin=1.5em]
    \item Acumulación de polvo que reduce el brillo de las imágenes.
    \item Condensación que favorece corrosión y crecimiento fúngico.
    \item Hongos ópticos que degradan irreversiblemente la transmitancia del vidrio.
    \item Entrada de insectos o cuerpos extraños al interior del tubo.
\end{itemize}

\begin{peligro}{Una tarde sin tapas en temporada de lluvias}{pel:hongos}
    \itshape
    Isabella termina una sesión de observación terrestre y guarda rápidamente el
    telescopio en el depósito del colegio, olvidando colocar la tapa del objetivo.
    Esa tarde empieza una semana de lluvias intensas; el depósito sube al 85\% de
    humedad relativa. En el transcurso de tres semanas, esporas de \textit{Aspergillus}
    presentes en el aire se depositan sobre la lente correctora húmeda del Maksutov
    y germinan.

    Dos meses después, durante una sesión nocturna, el grupo nota que la imagen de
    la Luna tiene un velo difuso que no existía antes. Al revisar con una linterna,
    se descubren manchas blanquecinas irregulares sobre la lente: hongos ópticos.
    El telescopio debe enviarse a un técnico especializado. El diagnóstico: el
    recubrimiento antirreflejo de la lente correctora necesita tratamiento químico
    de profundidad y posiblemente reacondicionamiento.

    \medskip
    \normalfont
    El costo de las tapas es cero. El costo de ignorarlas puede ser la inutilización
    del instrumento durante meses. \cite{rutten1988,cloudy2019}
\end{peligro}

% ---------------------------------------------------------------
\section{Esteban — No forzar los mecanismos de ajuste del telescopio}
% ---------------------------------------------------------------

\begin{concepto}{Precisión mecánica y tolerancias en telescopios}{con:mecanica}
    Un telescopio es un instrumento de \emph{precisión mecánica}: sus mecanismos
    de enfoque, ajuste de montura y colimación trabajan con tolerancias de décimas
    de milímetro. El mecanismo helicoide de enfoque de un Maksutov, por ejemplo,
    mueve la lente correctora —~o el conjunto de espejos~— una distancia controlada
    por la rosca del helicoide, cuyo paso puede ser de apenas 0.5~mm por vuelta
    \cite{tonkin2007}. Cualquier fuerza excesiva puede deformar roscas, doblar
    el porta-oculares o, en el caso de la montura, desalinear el eje óptico de
    todo el sistema.
\end{concepto}

\subsection{¿Por qué es importante esta norma?}

Los mecanismos de ajuste pueden «resistirse» por varias razones completamente normales:

\begin{itemize}[leftmargin=1.5em]
    \item \textbf{Dilatación o contracción térmica}: el metal se expande y contrae
          al cambiar de temperatura; en la mañana fría una rosca apretada por calor
          puede parecer atascada. Solución: esperar la aclimatación.
    \item \textbf{Lubricante envejecido}: las grasas de las roscas se solidifican
          con el tiempo. Solución: mantenimiento preventivo con grasa de silicona.
    \item \textbf{Polvo o suciedad en la rosca}: partículas que crean fricción extra.
          Solución: limpieza con aire comprimido antes de insistir.
\end{itemize}

\begin{nota}{La colimación: el parámetro óptico más sensible}{nota:colimacion}
    La \emph{colimación} es la alineación precisa del eje óptico de todos los
    elementos del telescopio. En el Maksutov-Cassegrain, esta alineación es muy
    estable gracias al diseño cerrado, pero puede desajustarse si se aplica
    torsión o impacto a la lente correctora o al tubo \cite{suiter1994}. Un
    telescopio descolimado produce imágenes astigmáticas: las estrellas en el
    centro del campo aparecen como elipses o cometas, nunca como puntos.
\end{nota}

\subsection{¿Qué riesgo se evita?}

\begin{itemize}[leftmargin=1.5em]
    \item Daño al mecanismo de enfoque helicoide (rosca deformada o porta-oculares torcido).
    \item Descolimación de los espejos por esfuerzo torsional sobre el tubo.
    \item Daño a la montura alt-azimutal (ejes doblados o tornillos de tensión rotos).
    \item Inoperabilidad del telescopio, que requeriría reparación técnica especializada.
\end{itemize}

\begin{peligro}{El enfocador que nunca volvió a funcionar bien}{pel:forzar}
    \itshape
    En una sesión matutina, Esteban intenta enfocar la imagen de un edificio lejano.
    El helicoide de enfoque se siente rígido porque el telescopio todavía no se
    ha aclimatado y el metal está contraído por la temperatura de la madrugada.
    Esteban, pensando que está atascado por suciedad, aplica fuerza progresiva
    con ambas manos. En un momento, siente un «clic» y el enfocador se mueve.

    Lo que Esteban no sabe es que ese «clic» fue el stripped de la primera rosca
    del helicoide. A partir de ese día, el mecanismo de enfoque tiene juego axial:
    el tubo interior se tambalea ligeramente al girar, haciendo imposible mantener
    el enfoque preciso a altos aumentos. La imagen de Júpiter nunca estará del todo
    nítida. Reparar el helicoide de un Maksutov requiere desensamblar el tubo
    óptico completo, un trabajo de relojería que ningún aficionado puede hacer en
    el colegio.

    \medskip
    \normalfont
    Si algo no se mueve con suavidad, la respuesta correcta es investigar la causa,
    no aplicar más fuerza. \cite{tonkin2007,suiter1994}
\end{peligro}

% ---------------------------------------------------------------
\section{Samuel Angulo — Funcionamiento del sistema de aumento del telescopio}
% ---------------------------------------------------------------

\begin{concepto}{Magnificación en telescopios refractores y catadióptricos}{con:aumento}
    El \emph{aumento} (o magnificación) de un telescopio se define como la relación
    entre la distancia focal del instrumento y la del ocular utilizado:
    \[
        M \;=\; \frac{F_{\text{telescopio}}}{F_{\text{ocular}}}
    \]
    donde $F_{\text{telescopio}}$ es la distancia focal del telescopio (en mm) y
    $F_{\text{ocular}}$ es la distancia focal del ocular (en mm). Para el
    \emph{Decouvir Maksutov-Cassegrain MC~60/700}, la distancia focal es de
    aproximadamente 700~mm, con apertura de 60~mm y relación focal $f/11.7$
    \cite{decouvir2020}.
\end{concepto}

\subsection{El telescopio Maksutov-Cassegrain: diseño catadióptrico}

El \textit{Decouvir MC~60/700} es un telescopio \textbf{catadióptrico} de diseño
Maksutov-Cassegrain \cite{rutten1988}. A diferencia de un refractor (solo lentes)
o un newtoniano (solo espejos), combina ambos elementos:

\begin{itemize}[leftmargin=1.5em]
    \item \textbf{Lente correctora de menisco} (en la parte frontal del tubo):
          una lente cóncava-convexa que corrige las aberraciones del espejo primario
          esférico y actúa como «tapa» hermética del tubo.
    \item \textbf{Espejo primario esférico} (en el fondo del tubo): concentra la luz
          hacia el espejo secundario.
    \item \textbf{Espejo secundario convexo} (mancha plateada en el centro de la
          lente correctora): refleja la luz de vuelta por un agujero en el espejo
          primario hacia el ocular.
\end{itemize}

Esta disposición «pliega» el camino óptico dentro de un tubo muy compacto,
logrando una distancia focal de 700~mm en un tubo de apenas $\sim$15~cm de
longitud. El diseño cerrado también protege los espejos del polvo y la humedad
de forma inherente \cite{tonkin2007}.

\subsection{Oculares y aumentos disponibles}

Con los oculares más comunes para telescopios de iniciación, los aumentos
aproximados en el \textit{MC~60/700} serían:

\begin{center}
\begin{tabular}{lcc}
\hline
\textbf{Ocular (mm)} & \textbf{Aumento ($M$)} & \textbf{Uso recomendado} \\
\hline
25 & $700/25 = 28\times$ & Cúmulos, nebulosas, visión panorámica \\
20 & $700/20 = 35\times$ & Luna completa, campo amplio \\
12.5 & $700/12.5 = 56\times$ & Cráter lunar, planetas brillantes \\
10  & $700/10 = 70\times$ & Detalle lunar, bandas de Júpiter \\
4   & $700/4 = 175\times$ & Límite práctico para 60~mm \\
\hline
\end{tabular}
\end{center}

\begin{regla}{Aumento máximo útil para 60~mm de apertura}{reg:max-aum}
    El aumento máximo \emph{útil} de cualquier telescopio se estima como
    $M_{\max} \approx 2 \times D$ (donde $D$ es la apertura en milímetros).
    Para 60~mm de apertura:
    \[
        M_{\max} \approx 2 \times 60 = 120\times
    \]
    Por encima de este valor, las imágenes se vuelven más grandes pero más
    difusas —~lo que se llama \emph{magnificación vacía}~— sin revelar más
    detalle real \cite{suiter1994,rutten1988}. En condiciones de turbulencia
    atmosférica (\emph{seeing} deficiente), incluso $70\times$ puede ser
    excesivo.
\end{regla}

\subsection{¿Qué riesgo se evita al entender el sistema de aumento?}

\begin{itemize}[leftmargin=1.5em]
    \item \textbf{Frustración y desmotivación} al usar aumentos inadecuados para
          las condiciones atmosféricas.
    \item \textbf{Uso incorrecto de la lente de Barlow}: una Barlow $2\times$
          multiplica el aumento del ocular por dos, lo que puede llevar al aumento
          vacío si no se calcula correctamente.
    \item \textbf{Elección incorrecta del ocular} para el objeto de interés: no
          todos los objetos necesitan el máximo aumento.
\end{itemize}

\begin{idea}{La regla del «aumento progresivo»}{idea:progresivo}
    Una buena práctica para cualquier observación es comenzar siempre con el
    ocular de \emph{menor aumento} (mayor distancia focal), centrar el objeto
    en el campo visual, y luego aumentar gradualmente hasta encontrar el aumento
    óptimo para las condiciones de esa noche. Nunca empezar con el máximo aumento
    \cite{tonkin2007}.
\end{idea}

\begin{peligro}{Saturno a 175× en una noche de \emph{bad seeing}}{pel:aumento}
    \itshape
    Samuel Angulo, emocionado por ver Saturno por primera vez, coloca directamente
    el ocular de 4~mm (175×) en el telescopio. La imagen que aparece es grande,
    sí, pero también borrosa, inestable y decepcionante: los anillos se ven como
    un blob difuso, sin definición. Samuel concluye que el telescopio es de mala
    calidad y pierde interés.

    Lo que Samuel no sabe es que esa noche el índice de \emph{seeing} es de apenas
    2/5 por la turbulencia atmosférica. Con el ocular de 20~mm (35×), la imagen de
    Saturno habría sido pequeña pero perfectamente nítida, con los anillos claramente
    separados del disco y la división de Cassini casi visible. A 70× también habría
    sido excelente. Fue el aumento excesivo, no el telescopio, el que arruinó la
    observación.

    \medskip
    \normalfont
    Comprender el sistema de aumento permite calibrar las expectativas y extraer
    lo mejor del instrumento en cada condición atmosférica. \cite{suiter1994,tonkin2007}
\end{peligro}

% ---------------------------------------------------------------
\section{Síntesis: las normas como sistema integrado}
% ---------------------------------------------------------------

Las seis normas asignadas no son una lista de prohibiciones aisladas: forman un
\textbf{sistema integrado de buenas prácticas} cuyo objetivo es preservar el
instrumento y maximizar la calidad de cada sesión de observación.

\begin{actividad}{Principio unificador de las normas de uso}{act:sintesis}
    Cada norma responde a uno de tres principios fundamentales:
    \begin{enumerate}[leftmargin=1.5em]
        \item \textbf{Seguridad personal} (norma de Miguel): proteger la integridad
              física de los observadores.
        \item \textbf{Preservación del instrumento} (normas de Kevin, Isabella, Esteban):
              mantener el telescopio en condiciones óptimas para el máximo tiempo posible.
        \item \textbf{Calidad de la observación} (normas de Thomas y Samuel): garantizar
              que cada sesión produzca resultados útiles y satisfactorios.
    \end{enumerate}
    Un semillero que internalice estos tres principios estará preparado no solo
    para usar el \textit{Decouvir MC~60/700}, sino cualquier instrumento óptico
    que encuentre en su camino.
\end{actividad}

\begin{compromiso}{Compromisos del semillero con el instrumento}{comp:cuidado}
    \begin{itemize}[leftmargin=1.5em]
        \item Nunca apuntar al Sol sin filtro solar certificado ISO~12312-2.
        \item Aclimatar el telescopio al menos 20 minutos antes de observar.
        \item Nunca tocar superficies ópticas con los dedos ni con telas no apropiadas.
        \item Colocar las tapas del objetivo y del ocular al guardar el instrumento.
        \item Investigar siempre la causa de una resistencia mecánica antes de forzar.
        \item Usar siempre el ocular de menor aumento primero y aumentar gradualmente.
    \end{itemize}
\end{compromiso}

% ---------- BIBLIOGRAFÍA ----------
\backmatter
\chapter*{Referencias}
\addcontentsline{toc}{chapter}{Referencias}

\begin{thebibliography}{99}

\bibitem{suiter1994}
Suiter, H.~R. (1994).
\textit{Star Testing Astronomical Telescopes: A Manual for Optical Evaluation and Adjustment}.
Willmann-Bell, Richmond, VA.

\bibitem{rutten1988}
Rutten, R.~A.~M., y van Venrooij, M.~A.~M. (1988).
\textit{Telescope Optics: Evaluation and Design}.
Willmann-Bell, Richmond, VA.

\bibitem{thompson2005}
Thompson, R., y Thompson, B. (2005).
\textit{Astronomy Hacks: Tips \& Tools for Observing the Night Sky}.
O'Reilly Media, Sebastopol, CA.

\bibitem{tonkin2007}
Tonkin, S.~F. (2007).
\textit{Practical Amateur Spectroscopy}.
Springer, Londres.

\bibitem{pasachoff2000}
Pasachoff, J.~M. (2000).
\textit{A Field Guide to the Stars and Planets} (4.ª ed.).
Houghton Mifflin, Boston, MA.

\bibitem{smith2000}
Smith, W.~J. (2000).
\textit{Modern Optical Engineering: The Design of Optical Systems} (3.ª ed.).
McGraw-Hill, Nueva York.

\bibitem{dobson2011}
Dobson, R., y Bhatt, M. (2011).
Solar retinopathy in amateur astronomers.
\textit{British Journal of Ophthalmology}, 95(3), 427–428.
\url{https://doi.org/10.1136/bjo.2010.185330}

\bibitem{iso12312}
International Organization for Standardization. (2015).
\textit{ISO 12312-2:2015 — Eye and face protection. Sunglasses and related eyewear.
Part 2: Filters for direct observation of the Sun}.
ISO, Ginebra.

\bibitem{decouvir2020}
Decouvir Optics. (2020).
\textit{Manual de usuario: Telescopio Maksutov-Cassegrain MC~60/700}.
Decouvir Optics.

\bibitem{baader2023}
Baader Planetarium GmbH. (2023).
\textit{AstroSolar Safety Film — Technical Data Sheet}.
Baader Planetarium, Mammendorf.
Recuperado de \url{https://www.baader-planetarium.com/en/baader-astrosolar-film.html}

\bibitem{cloudy2019}
Cloudy Nights Community Forums. (2019).
\textit{Fungus on telescope optics: prevention and treatment}.
Recuperado de \url{https://www.cloudynights.com/topic/fungus-on-optics}

\end{thebibliography}

\end{document}
